% !TeX TS-program = pdflatex


\documentclass[a4paper]{article}

% \usepackage[default]{fontsetup}

\usepackage{fancyhdr}
\usepackage{extramarks}
\usepackage{amsmath}
\usepackage{amsthm}
\usepackage{amsfonts}
\usepackage{tikz}
\usepackage[plain]{algorithm}
\usepackage{algpseudocode}
\usepackage{enumerate}
\usepackage{tikz}

\usetikzlibrary{automata,positioning}

%
% Basic Document Settings
%  

\topmargin=-0.2in
\evensidemargin=0in
\oddsidemargin=0in
\textwidth=6.5in
\textheight=9.5in
\headsep=0.25in

\linespread{1.1}

\pagestyle{fancy}
\lhead{\hmwkAuthorName}
\chead{\hmwkClass : \hmwkTitle}
\rhead{\firstxmark}
\lfoot{\lastxmark}
\cfoot{\thepage}

\renewcommand\headrulewidth{0.4pt}
\renewcommand\footrulewidth{0.4pt}

\setlength\parindent{0pt}

%
% Create Problem Sections
%

\newcommand{\enterProblemHeader}[1]{
    \nobreak\extramarks{}{Problem \arabic{#1} continued on next page\ldots}\nobreak{}
    \nobreak\extramarks{Problem \arabic{#1} (continued)}{Problem \arabic{#1} continued on next page\ldots}\nobreak{}
}

\newcommand{\exitProblemHeader}[1]{
    \nobreak\extramarks{Problem \arabic{#1} (continued)}{Problem \arabic{#1} continued on next page\ldots}\nobreak{}
    \stepcounter{#1}
    \nobreak\extramarks{Problem \arabic{#1}}{}\nobreak{}
}

\newcommand*\circled[1]{\tikz[baseline=(char.base)]{
		\node[shape=circle,draw,inner sep=2pt] (char) {#1};}}


\setcounter{secnumdepth}{0}
\newcounter{partCounter}
\newcounter{homeworkProblemCounter}
\setcounter{homeworkProblemCounter}{1}
\nobreak\extramarks{Problem \arabic{homeworkProblemCounter}}{}\nobreak{}

%
% Homework Problem Environment
%
% This environment takes an optional argument. When given, it will adjust the
% problem counter. This is useful for when the problems given for your
% assignment aren't sequential. See the last 3 problems of this template for an
% example.
%

\newenvironment{homeworkProblem}[1][-1]{
    \ifnum#1>0
        \setcounter{homeworkProblemCounter}{#1}
    \fi
    \section{Problem \arabic{homeworkProblemCounter}}
    \setcounter{partCounter}{1}
    \enterProblemHeader{homeworkProblemCounter}
}{
    \exitProblemHeader{homeworkProblemCounter}
}

%
% Homework Details
%   - Title
%   - Class
%   - Due date
%   - Name
%   - Student ID

\newcommand{\hmwkTitle}{Homework\ \#07}
\newcommand{\hmwkClass}{Probability \& Statistics for EECS}
\newcommand{\hmwkDueDate}{2026-11-24}
\newcommand{\hmwkAuthorName}{}
\newcommand{\hmwkAuthorID}{}


%
% Title Page
%

\title{
    \vspace{2in}
    \textmd{\textbf{\hmwkClass:\\  \hmwkTitle}}\\
    \normalsize\vspace{0.1in}\small{Due\ on\ Nov. 10, 2025, 23:59 UTC+8}\\
	\vspace{4in}
}

\author{
	Name: \textbf{\hmwkAuthorName} \\
	Student ID: \hmwkAuthorID}
\date{}

\renewcommand{\part}[1]{\textbf{\large Part \Alph{partCounter}}\stepcounter{partCounter}\\}

%
% Various Helper Commands
%

% Useful for algorithms
\newcommand{\alg}[1]{\textsc{\bfseries \footnotesize #1}}
% For derivatives
\newcommand{\deriv}[1]{\frac{\mathrm{d}}{\mathrm{d}x} (#1)}
% For partial derivatives
\newcommand{\pderiv}[2]{\frac{\partial}{\partial #1} (#2)}
% Integral dx
\newcommand{\dx}{\mathrm{d}x}
% Alias for the Solution section header
\newcommand{\solution}{\textbf{\large Solution}}
% Probability commands: Expectation, Variance, Covariance, Bias
\newcommand{\E}{\mathrm{E}}
\newcommand{\Var}{\mathrm{Var}}
\newcommand{\Cov}{\mathrm{Cov}}
\newcommand{\Bias}{\mathrm{Bias}}

\begin{document}


% \maketitle
% \thispagestyle{empty}
% \pagebreak

\date{
Due on Nov. 24, 2026, 18:00 UTC+8}
\title{SI 140A-01  Probability \& Statistics for EECS, Fall 2026 \\
Homework 7}
\maketitle
Read all the instructions below carefully before you start working on the assignment, and before you make a submission.
\begin{itemize}
    \item You are required to write down all the major steps towards making your conclusions; otherwise you may obtain limited points of the problem.
    \item Write your homework in English; otherwise you will get no points of this homework.
    \item Any form of plagiarism will lead to $0$ point of this homework. 
\end{itemize}
\newpage

\begin{homeworkProblem}[1] \textbf{(30pts)}
Let $U_1,\dots,U_n$ be i.i.d. $\text{Unif}(0,1)$, and $X = \max(U_1,\dots,U_n)$. What is the PDF of $X$? What is $E X$?

\smallskip

Hint: Find the CDF of $X$ first, by translating the event $X \leq x$ into an event involving $U_1,\dots,U_n$
\end{homeworkProblem}
\newpage

\begin{homeworkProblem}[2] \textbf{(30pts)}
Let $Z \sim \mathcal{N}(0,1)$ and let $S$ be a random sign independent of $Z$, i.e., $S$ is $1$ with probability $1/2$ and $-1$ with probability $1/2$. Show that $SZ \sim \mathcal{N}(0,1)$.
\end{homeworkProblem}

\newpage
\begin{homeworkProblem}[3] \textbf{(40pts)}
The \emph{Gumbel distribution} is the distribution of $-\log X$ with $X \sim \text{Expo}(1)$.

\begin{enumerate}
    \item Find the CDF of the Gumbel distribution.
    \item Let $X_1, X_2, \dots$ be i.i.d. $\text{Expo}(1)$ and let $M_n = \max(X_1, \dots, X_n)$. Show that $M_n - \log n$ converges in distribution to the Gumbel distribution, i.e., as $n \to \infty$ the CDF of $M_n - \log n$ converges to the Gumbel CDF.
\end{enumerate}
\end{homeworkProblem}

\newpage
\begin{homeworkProblem}[4]\textbf{(optional)}
Let $Y \sim \mathcal{N}(\mu, 1)$. Use the fact that $P(|Y - \mu| < 1.96\sigma) \approx 0.95$ to construct a random interval $(a(Y), b(Y))$ (that is, an interval whose endpoints are r.v.s), such that the probability that $\mu$ is in the interval is approximately $0.95$. This interval is called a \emph{confidence interval} for $\mu$; such intervals are often desired in statistics when estimating unknown parameters based on data.
\end{homeworkProblem}
 
 \newpage
 \begin{homeworkProblem}[5]
    Let $F$ be a CDF which is continuous and strictly increasing.
    Let $\mu$ be the mean of the distribution.
    The quantile function, $F^{-1}$, has many applications in statistics and econometrics. 
    Show that the area under the curve of the quantile function from $0$ to $1$ is $\mu$.\\

  

\end{homeworkProblem}



\end{document}